\section{Background}
\label{sec:background}

\subsection{Urban Radio Propagation}

Urban radio propagation at sub-6\,GHz frequencies in outdoor environments differs
fundamentally from purely outdoor line-of-sight propagation owing to the dense
arrangement of buildings and other scatterers that give rise to a rich multipath
environment~\cite{rappaport2013millimeter}.
The dominant propagation mechanisms in urban micro-cells are specular reflection
from building facades, diffraction at building edges and corners, and transmission
through exterior walls, each of which is governed by the frequency-dependent
electromagnetic properties of the building materials present~\cite{itu_p2040}.
These effects collectively determine the key channel descriptors used in system
design---most notably the root-mean-square (RMS) delay spread, the angular spread of
arrival and departure, and the aggregate path loss~\cite{3gpp_nr_channel}.
Accurate characterisation of such parameters is a prerequisite for designing link
budgets, dimensioning cyclic-prefix lengths in OFDM transmission, and exploiting
spatial multiplexing with multi-antenna systems~\cite{dahlman2018_5gnr}.

\subsection{Ray-Tracing Methods for Channel Modelling}

Ray-tracing (RT) is a deterministic propagation modelling technique that predicts
multipath channel characteristics by simulating the geometric and electromagnetic
interactions of a set of rays launched from a transmitting antenna with the surfaces
of a three-dimensional scene~\cite{yun2015ray}.
Compared with stochastic geometry-based models such as those specified in
3GPP~TR~38.901~\cite{3gpp_nr_channel}, RT methods directly exploit site-specific
scene and material information, making them particularly well-suited to the
reproduction and analysis of physically realistic channel propagation in urban environments.

\textbf{MATLAB Ray Tracing Toolbox} is employed in this work to generate the wireless
channel dataset using geometric optics ray tracing with configurable path-tracing parameters
(reflections, diffractions, frequency). The ray-tracing simulation is applied to an
OpenStreetMap-based three-dimensional model of the urban Otaniemi campus with material
properties assigned to building components, producing physically realistic multipath
channel characteristics for both localization and channel-charting applications.

